\section{Enlargeability}
\label{sec:enl}

In this section we discuss the concept of enlargeability, which is the central idea the results in this thesis are concerned with. 

\begin{definition}
    \label{enl:def:enlargeability}
    A closed oriented Riemannian $n$-manifold $M$ is called enlargeable if for all $\epsilon >0$ there exists a Riemannian covering $\hat{M} \to M$ together with an $\epsilon$-contracting (that is $\epsilon$-Lipschitz) map $\hat{M} \to S^n$, which is constant outside a compact subset and of non-zero degree.
\end{definition}

Note that this notion does not depend on the particular choice of Riemannian metric on $M$, since for closed manifolds all metrics are in bi-Lipschitz correspondence. Intuitively, the covering $\hat{M}$ can be pictured as an ``expansion'' of $M$ itself (think for example of the usual sketches of coverings of $S^1$). Enlargeability now means that this expansion can be chosen arbitrarily large, such that it can be wrapped as tightly around a sphere as one wishes. Naively one might expect that for an enlargeable manifold $M$ the universal cover $\tilde{M}$ - which in some sense is the ``largest'' covering - is hyperspherical, meaning that it admits an $\epsilon$-contracting map $\tilde{M} \to S^n$, which is constant outside a compact subset and of non-zero degree for arbitrarily small $\epsilon > 0$. However, Brunnbauer and Hanke have constructed counterexamples to this naive expectation in [BH, Theorem 1.4].

The simplest example of an enlargeable manifold is the Torus $T^n$, since it is covered by $\R^n$, which is clearly hyperspherical. More generally the same holds for any closed manifold admitting a metric of non-positive sectional curvature by the Cartan-Hadamard Theorem. Non-enlargeable manifolds are for example all simply connected manifolds, since no interesting covers exist in this case.

% Our definition of enlargeability is adopted from \cite{cecchini2021enlargeable} and \cite{HankeSchick} in contrast to the orginial definition in \cite{gromov1980spin}, where the coverings $\hat{M}$ were additionally required to be spin. 

% The simplest example of an enlargeable manifold is $S^1$. To see this consider the Riemannian covering $(S^1,n g) \to (S^1, g), z \mapsto z^n$, where $g$ is the standard metric on $S^1$. Then the map $(S^1,n g) \to (S^1, g), z \mapsto z$ is $\sfrac{1}{n}$-Lipschitz and of degree $1$. Since we can choose $n$ arbitrarily large it follows that $S^1$ is enlargeable. 

% PICTURE (usual covering)

% The observation that enlargeability is an obstruction to the existence of positive scalar curvature metrics was first shown by Gromov and Lawson \cite{gromov1980spin} in the case where $M$ is spin. This was later generalized by Cecchini and Schick \cite{cecchini2021enlargeable} to the non-spin case.

The central fact about enlargeable manifolds is that they do not admit psc metrics.

\begin{theorem}[GL1980, GL1982, CS2020]
    Enlargeable manifolds do not admit metrics of positive scalar curvature.
\end{theorem}

This is the result that motivated the original definition by Gromov and Lawson [GL1980], where they first showed that enlargeability is an obstruction to the existence of psc metrics, under the additional assumption that the coverings $\hat{M}$ in Definition \ref{enl:def:enlargeability} are finite an spin. They were later able to drop the finiteness assumption [GL1982]. Both proofs are based on techniques using the spin Dirac operator, so it was unclear how and if these results carry over to the non-spin case. Rather recently however, Cecchini and Schick resolved this issue in [CS2020] by relying on modern minimal surface methods developed by Schoen and Yau [SY...].

The following Lemma now provides two easy ways to construct new enlargeable manifolds out of given ones. 

\begin{lemma}
    \label{enl:lem:constructing_enlargeable_mnf}
    \begin{enumerate}[label=(\roman*)]
        \item The product of enlargeable manifolds is enlargeable. \label{enl:eq:lem_for_constructing_enlargeable_mnf_1}
        \item A closed oriented manifold admitting a non-zero degree map onto an enlargeable manifold is enlargeable. \label{enl:eq:lem_for_constructing_enlargeable_mnf_2}
    \end{enumerate}
\end{lemma}

\begin{proof}
    For \ref{eq:lem_for_constructing_enlargeable_mnf_1} let $M^m$ and $N^n$ be enlargeable. For $\epsilon > 0$ choose Riemannian coverings $p \colon \hat{M} \to M$ and $q \colon \hat{N} \to N$ as well as $\epsilon$-contracting maps $f \colon \hat{M} \to S^m$ and $g \colon \hat{N} \to S^n$ of non-zero degree. Then 
    \[
    p \times q \colon \hat{M} \times \hat{N} \to M \times N
    \]
    is a Riemannian covering, where both manifolds are equipped with the corresponding product metric. Moreover 
    \[
    f \times g \colon \hat{M} \times \hat{N} \to S^m \times S^n
    \]
    defines an $\epsilon$-contracting map of degree $\mathrm{deg}(f \times g) = \mathrm{deg}(f)\mathrm{deg}(g) \neq 0$. Now consider the map 
    \[
    \phi \colon S^m \times S^n \to S^m \wedge S^n \cong S^{m+n},
    \]
    given by projection onto the smash product 
    \[
    S^m \wedge S^n := \raisebox{3pt}{$\modulo{S^m \times S^n}{(S^m \times \ast \cup \ast \times S^n)}$},
    \]
    which is well known to be homeomorphic to the sphere $S^{m+n}$. The map $\phi$ is easily seen to be of degree $\pm 1$ (depending on the chosen orientations) and after a slight perturbation we may also assume it to be smooth. Moreover, because $S^m \times S^n$ is closed, it follows that $\phi$ is a $\lambda$-Lipschitz map for some constant $\lambda > 0$. Therefore $\phi \circ (f \times g)$ is a $\lambda \epsilon$-Lipschitz map of non-zero degree and since $\epsilon$ was chosen arbitrarily the statement follows.

    For \ref{eq:lem_for_constructing_enlargeable_mnf_2} let $M^m$ be enlargeable and let $h \colon P \to M$ be some non-zero degree map. Again given $\epsilon > 0$ choose a Riemannian covering $p \colon \hat{M} \to M$ as well as an $\epsilon$-contracting map $f \colon \hat{M} \to S^m$. Define a covering $\pi \colon \hat{P} \to P$ as the pullback covering of $p$ under $h$, that is $\pi$ fits into a pullback diagram of the form 
    \[\begin{tikzcd}
	{\hat{P}} & {\hat{M}} \\
	P & M
	\arrow["H", from=1-1, to=1-2]
	\arrow["\pi", from=1-1, to=2-1]
	\arrow["p", from=1-2, to=2-2]
	\arrow["h", from=2-1, to=2-2]
    \end{tikzcd}\]
    where $H$ is some suitable smooth map. Since $\hat{M}$ and $P$ are closed, the same holds for $\hat{P}$. Moreover $h$ is $\mu$-Lipschitz for some constant $\mu > 0$, because $P$ is closed. Consequently $H$ is also $\mu$-Lipschitz, given that both $\pi$ and $p$ are local isometries. For the degree of $H$ we compute 
    \[
    \mathrm{deg}(p)\mathrm{deg}(H) = \mathrm{deg}(p \circ H) = \mathrm{deg}(h \circ \pi) = \mathrm{deg}(h)\mathrm{deg}(\pi) \neq 0,
    \]
    using the fact that $\mathrm{deg}(h) \neq 0$ by assumption and that the degree of a covering is given by the cardinality of the fiber. We conclude that $H$ is of non-zero degree, such that $f \circ H \colon \hat{P} \to S^m$ is a $\mu \epsilon$-Lipschitz map of non-zero degree. Again since $\epsilon$ was chosen arbitrarily the statement follows.
\end{proof}

From Part \ref{eq:lem_for_constructing_enlargeable_mnf_2} of the above lemma it follows immediately, that enlargeability is a homotopy invariant, since a homotopy equivalence of connected closed oriented manifolds has degree $\pm 1$. In fact much more is true, namely that enlargeability is a homological invariant in the following sense. 

\begin{theorem*}[\cite{BrunnbauerHanke}]
    Let $M$ be a closed oriented $n$-manifold and let $\phi \colon M \to B\pi_1(M)$ be the classifying map of the universal cover. Then $M$ is enlargeable if and only if $\phi_\ast [M] \in H_n(B\pi_1(M);\Q)$ does not lie in the rational vector subspace 
    \[
    H_n^\mathrm{small}(B\pi_1(M);\Q) \subseteq H_n(B\pi_1(M);\Q)
    \]
    of ``small'' classes, i.e. $\phi_\ast [M]$ is ``large''.
    % In fact the assignment $M \mapsto H_n^\mathrm{small}(B\pi_1(M);\Q)$ is functorial for maps $f \colon M \to N$ inducing an epimorphism on $\pi_1$, that is $f_\ast \colon H_n(B\pi_1(M);\Q) \to H_n(B\pi_1(N);\Q)$ maps small classes to small classes.
\end{theorem*}

 In \cite{BrunnbauerHanke} the authors showed this by giving a general definition of ``small'' and ``large'' homology classes of a simplicial complex $X$ with finitely generated fundamental group, such that the assignment $X \mapsto H_n^\mathrm{small}(X;\Q)$ is functorial for continuous maps $f \colon X \to Y$ inducing an epimorphism on $\pi_1$, that is $f_\ast \colon H_n(X;\Q) \to H_n(Y;\Q)$ maps small classes to small classes. Moreover if $f$ happen to induce an isomorphism on $\pi_1$ then $f_\ast$ also maps large classes to large classes. This general definition then specializes to the fact that a closed oriented manifold $M$ is enlargeable if and only if the fundamental class $[M]$ is large, that is $[M] \notin H_n^\mathrm{small}(M;\Q)$. Combining these observations, the theorem follows. Furthermore this argument shows that the notion of enlargeability can be purely thought of as a concept of simplicial complexes. For details we refer the reader to \cite{BrunnbauerHanke}.

% \begin{theorem*}[\cite{BrunnbauerHanke}]
%     Let $M$ be a closed oriented $n$-manifold and let $\phi \colon M \to B\pi_1(M)$ be the classifying map of the universal cover. Then there exists a rational vector subspace 
%     \[
%     H_n^\mathrm{small}(B\pi_1(M);\Q) \subseteq H_n(B\pi_1(M);\Q)
%     \]
%     of ``small'' classes with the property that $M$ is enlargeable if and only if $\phi_\ast [M] \notin H_n^\mathrm{small}(B\pi_1(M);\Q)$, that is $\phi_\ast [M]$ is ``large''.
%     % In fact the assignment $M \mapsto H_n^\mathrm{small}(B\pi_1(M);\Q)$ is functorial for maps $f \colon M \to N$ inducing an epimorphism on $\pi_1$, that is $f_\ast \colon H_n(B\pi_1(M);\Q) \to H_n(B\pi_1(N);\Q)$ maps small classes to small classes.
% \end{theorem*}



% In particular the above Theorem implies homotopy invariance of enlargeability. In \cite{BrunnbauerHanke} the authors showed this by giving a general definition of ``small'' and ``large'' homology classes of a simplicial complex $X$ with finitely generated fundamental group, such that the assignment $X \mapsto H_n^\mathrm{small}(X;\Q)$ is again functorial for continuous maps inducing an epimorphism on $\pi_1$. This definition then specializes to the fact that a closed oriented manifold $M$ is enlargeable if and only if the fundamental class $[M]$ is large, that is $[M] \notin H_n^\mathrm{small}(M;\Q)$, from which the Theorem follows. Moreover this argument shows that the notion of enlargeability can be purely thought of as a concept of simplicial complexes. For details we refer the reader to \cite{BrunnbauerHanke}.
